% Propietario conservado: /Users/ruben/Documents/New project/output/EXTREMA_MEDIA_RAZON_RECIPROCIDAD_ES_EN_20260918/ARTICULO/ES/source/nuclear/07.tex
\chapter{Variational conservation and oriented action}
\label{x_reciprocidad:nuclear:7}
\section{Symmetry of the elliptic action and Noether charge}

Fix a positive action section \(\hbar\), the angular pair \(A,C^*\) in a common unit, \(A\ne0\), \(|C^*/A|<1\), and

\[
\eta=\operatorname{artanh}(C^*/A).
\]

The realized plane \(\mathcal P_\eta\) has coordinates \((Q,P)\), both with the dimension of the square root of action, with

\[
\lambda=P\,dQ,\qquad \Omega=dQ\wedge dP=-d\lambda.
\]

The construction in \cref{x_reciprocidad:iv:ellipse:section} gives

\begin{equation}
\mathcal I_\eta(Q,P)=\frac12(e^{-\eta}Q^2+e^\eta P^2),
\qquad H_\eta=\omega\mathcal I_\eta,\qquad \omega>0.
\label{x_reciprocidad:nuc07:eq:1}
\end{equation}

The HMT ellipse is \(\mathcal I_\eta=\hbar\). Its area is \(2\pi\hbar\); \(\mathcal I_\eta\) has the dimension of action and \(H_\eta\) that of energy. Define the already realized matrices

\begin{equation}
M_\eta=\operatorname{diag}(e^{\eta/2},e^{-\eta/2}),\quad
J_2=\begin{pmatrix}0&-1\\1&0\end{pmatrix},\quad
R_\eta(\varepsilon)=M_\eta e^{-\varepsilon J_2}M_\eta^{-1}.
\label{x_reciprocidad:nuc07:eq:2}
\end{equation}

\begin{proposition}
The group \(R_\eta\) is a canonical symmetry of the action

\begin{equation}
\mathscr S_\eta[z]=\int_{t_a}^{t_b}
\left(P\dot Q-\omega\mathcal I_\eta(Q,P)\right)dt
\label{x_reciprocidad:nuc07:eq:3}
\end{equation}

and its Noether charge is \(\mathcal I_\eta\).
\end{proposition}

\begin{proof}
The determinant of \(M_\eta\) is one, and \(M_\eta^{\mathsf T}J_2M_\eta=J_2\). Moreover, \(\mathcal I_\eta(M_\eta y)=\|y\|^2/2\), which is invariant under \(e^{-\varepsilon J_2}\).
The generator

\[
X_\eta=e^\eta P\,\partial_Q-e^{-\eta}Q\,\partial_P
\]

satisfies \(\iota_{X_\eta}\Omega=d\mathcal I_\eta\). By Cartan's formula,

\[
\mathcal L_{X_\eta}\lambda=dF_\eta,\qquad
F_\eta=\lambda(X_\eta)-\mathcal I_\eta=e^\eta P^2-\mathcal I_\eta.
\]

Since \(X_\eta H_\eta=0\), the variation of the action density is \(dF_\eta/dt\). Its charge is

\[
\frac{\partial L}{\partial\dot Q}X_\eta Q-F_\eta
=Pe^\eta P-F_\eta=\mathcal I_\eta.
\]

The equations \(\dot Q=\omega e^\eta P\), \(\dot P=-\omega e^{-\eta}Q\) directly verify \(d\mathcal I_\eta/dt=0\).
\end{proof}

The positive parametrization of the ellipse is \(Q=a_S\cos\theta\), \(P=b_S\sin\theta\). These equations assign it \(\dot\theta=-\omega\). Thus, over a turn traversed according to the dynamical flow, \(\oint P\,dQ=2\pi\mathcal I_\eta\); with increasing \(\theta\), the integral is \(-2\pi\mathcal I_\eta\). This is the orientation distinction established in \cref{x_reciprocidad:iv:ellipse:area}. On the radional contour, the magnitude is \(h=2\pi\hbar\).

\section{Change of form: determined transport and connection}

For two already published forms \(\eta,\eta'\) and a phase advance \(\vartheta\), define

\begin{equation}
\mathcal T_{\eta',\eta}^{\vartheta}:\mathcal P_\eta\to\mathcal P_{\eta'},
\qquad
\mathcal T_{\eta',\eta}^{\vartheta}
=M_{\eta'}e^{-\vartheta J_2}M_\eta^{-1}.
\label{x_reciprocidad:nuc07:eq:4}
\end{equation}

\begin{proposition}
This transport is symplectic and satisfies

\begin{equation}
\mathcal I_{\eta'}(\mathcal T_{\eta',\eta}^{\vartheta}z)
=\mathcal I_\eta(z),\qquad
\mathcal T_{\eta'',\eta'}^{\vartheta_2}
\mathcal T_{\eta',\eta}^{\vartheta_1}
=\mathcal T_{\eta'',\eta}^{\vartheta_1+\vartheta_2}.
\label{x_reciprocidad:nuc07:eq:5}
\end{equation}
\end{proposition}

\begin{proof}
All factors are symplectic. The first equality reduces to invariance of \(\|y\|^2/2\) under rotation. In the second, the adjacent matrices \(M_{\eta'}^{-1}M_{\eta'}\) cancel and the phases add.
\end{proof}

The formula receives the forms and advances from the reader; it does not select their values. The complete list of transitions can be retained alongside the composite transport: two lists yielding the same composite map are not identified as histories.

\begin{proposition}
A differentiable realization \(\eta(t),\vartheta(t)\in C^1\), with \(\dot\vartheta=\omega(t)\), of the transport \eqref{x_reciprocidad:nuc07:eq:4} has generator

\begin{equation}
\boxed{
H_{\rm tr}(Q,P,t)=
\omega(t)\mathcal I_{\eta(t)}(Q,P)
+\frac{\dot\eta(t)}2 QP.}
\label{x_reciprocidad:nuc07:eq:6}
\end{equation}

This generator preserves \(\mathcal I_{\eta(t)}(z(t))\).
\end{proposition}

\begin{proof}
Differentiating \(z(t)=M_{\eta(t)}e^{-\vartheta(t)J_2}y_0\) gives

\[
\dot Q=\omega e^\eta P+\frac{\dot\eta}2Q,\qquad
\dot P=-\omega e^{-\eta}Q-\frac{\dot\eta}2P.
\]

These are Hamilton's equations for \eqref{x_reciprocidad:nuc07:eq:6}. The connection term realizes exactly \(\dot M_\eta M_\eta^{-1}z\), and its Hamiltonian is determined up to a function of \(t\) alone. Indeed,

\begin{equation}
\begin{aligned}
\frac{d\mathcal I_\eta}{dt}
&=\frac{\dot\eta}{2}(-e^{-\eta}Q^2+e^\eta P^2)
+\left\{\mathcal I_\eta,\frac{\dot\eta}{2}QP\right\}\\
&=\frac{\dot\eta}{2}(-e^{-\eta}Q^2+e^\eta P^2)
+\frac{\dot\eta}{2}(e^{-\eta}Q^2-e^\eta P^2)=0.
\end{aligned}
\label{x_reciprocidad:nuc07:eq:7}
\end{equation}

Equivalently, \(Q=e^{\eta/2}q\), \(P=e^{-\eta/2}p\) give

\begin{equation}
P\dot Q-H_{\rm tr}
=p\dot q-\frac{\omega(t)}2(q^2+p^2).
\label{x_reciprocidad:nuc07:eq:8}
\end{equation}

The rotational symmetry and its charge are transported exactly.
\end{proof}

Here the connection term has been derived from transport, not from an external force. The fixed-form LC Hamiltonian preserves \(\mathcal I_\eta\). During a deformation, omitting the second term of \eqref{x_reciprocidad:nuc07:eq:6} leaves the first term of \eqref{x_reciprocidad:nuc07:eq:7} uncancelled. Variational conservation also requires transporting the canonical form.

The proposition determines the generator of the declared interpolation; it does not assign a unique physical interpolation to every TPK history. The discrete identity \eqref{x_reciprocidad:nuc07:eq:4}–\eqref{x_reciprocidad:nuc07:eq:5} works without interpolation.

\section{Orientation and boundary of the action}

The transformations in \cref{x_reciprocidad:iv:ellipse:reciprocidad} distinguish \(S(Q,P)=(P,Q)\) and \(J_2(Q,P)=(-P,Q)\). Both exchange the form \(\eta\leftrightarrow-\eta\), but

\begin{equation}
S^*\Omega=-\Omega,\quad J_2^*\Omega=\Omega,\qquad
S^*\lambda=-\lambda+d(PQ),\quad
J_2^*\lambda=\lambda-d(PQ).
\label{x_reciprocidad:nuc07:eq:9}
\end{equation}

Likewise,

\begin{equation}
S R_\eta(\varepsilon)S^{-1}=R_{-\eta}(-\varepsilon),\qquad
J_2R_\eta(\varepsilon)J_2^{-1}=R_{-\eta}(\varepsilon).
\label{x_reciprocidad:nuc07:eq:10}
\end{equation}

These are verified using \(SM_\eta=M_{-\eta}S\), \(J_2M_\eta=M_{-\eta}J_2\), \(SJ_2S^{-1}=-J_2\). For an oriented open curve \(z:[a,b]\to\mathcal P_\eta\),

\begin{equation}
\int_{J_2z}\lambda=\int_z\lambda-[PQ]_a^b,\qquad
\int_{Sz}\lambda=-\int_z\lambda+[PQ]_a^b.
\label{x_reciprocidad:nuc07:eq:11}
\end{equation}

The boundary terms vanish on a cycle and cancel when compatible endpoints are concatenated. Quadratic energy alone does not determine these signs or these terms.

Nonadic monodromy records phase return and \(q_{\rm mem}\mapsto q_{\rm mem}+1\). Constancy of a charge does not require the counter to remain fixed. In the elliptic realization, a turn preserves \(\mathcal I_\eta\) while increasing the accumulated oriented integral by \(2\pi\mathcal I_\eta\). Charge, integrated action, phase, and register have distinct types. An angular turn is not identified here with an arbitrary TPK cycle without specifying its reader.

\section{Variational clock charge and balance between reading and memory}

\subsection{An explicit action charge of the clock}

The calendar is realized on \(\mathcal H_{\rm clk}=\mathbb C^{108}\), with Fourier projectors \(P_k=|\widetilde k\rangle\langle\widetilde k|\), and determines

\begin{equation}
N_{108}=\sum_{k=0}^{107}kP_k,\quad
H_{\rm clk}=\hbar\omega_0N_{108},\quad
\omega_0=\frac{2\pi}{108t_0}.
\label{x_reciprocidad:nuc07:eq:12}
\end{equation}

Its propagator over \(t_0\) is the shift by one cell. The clock action is

\begin{equation}
\mathscr S[\psi]=\int
\left[\frac{i\hbar}{2}
(\langle\psi,\dot\psi\rangle-\langle\dot\psi,\psi\rangle)
-\langle\psi,H_{\rm clk}\psi\rangle\right]dt.
\label{x_reciprocidad:nuc07:eq:13}
\end{equation}

The inner product is antilinear in the first variable. The equation is \(i\hbar\dot\psi=H_{\rm clk}\psi\).

\begin{proposition}
For \(B=B^*\) with \([B,H_{\rm clk}]=0\), the symmetry \(\psi\mapsto e^{-i\varepsilon B}\psi\) has charge

\begin{equation}
\mathcal Q_B(\psi)=\hbar\langle\psi,B\psi\rangle.
\label{x_reciprocidad:nuc07:eq:14}
\end{equation}
\end{proposition}

\begin{proof}
Localizing the parameter, \(\delta\psi=-i\varepsilon(t)B\psi\), gives \(\delta L=\hbar\dot\varepsilon\langle\psi,B\psi\rangle\). Integration by parts in the variational equations gives \(d\mathcal Q_B/dt=0\). Directly,

\[
\frac{d}{dt}\langle\psi,B\psi\rangle
=\frac{i}{\hbar}\langle\psi,[H_{\rm clk},B]\psi\rangle=0.
\]

With \(B=I\), norm is conserved; with \(B=N_{108}\), the charge is the action \(E_{\rm clk}/\omega_0\).
\end{proof}

The symplectic form of the underlying real space is \(\Omega_{\rm clk}(u,v)=2\hbar\operatorname{Im}\langle u,v\rangle\); its one-form is \(\Theta_\psi(v)=-\hbar\operatorname{Im}\langle\psi,v\rangle\), with \(\Omega_{\rm clk}=-d\Theta\).

\subsection{Exact charge balance}

Let \(\mathcal H,\mathcal K\) be finite-dimensional Hermitian spaces, \(J:\mathcal H\to\mathcal K\) an isometry, \(U\) unitary on \(\mathcal K\), and

\begin{equation}
\Pi=JJ^*,\quad T=J^*UJ,\quad
\eta_{\rm mem}=(I-\Pi)UJ.
\label{x_reciprocidad:nuc07:eq:15}
\end{equation}

The operator \(\eta_{\rm mem}\) is distinct from the scalar parameter \(\eta\).

\begin{proposition}
If

\begin{equation}
\widetilde B=\widetilde B^*,\quad U^*\widetilde BU=\widetilde B,
\quad [\widetilde B,\Pi]=0,\quad B=J^*\widetilde BJ,
\label{x_reciprocidad:nuc07:eq:16}
\end{equation}

then

\begin{equation}
B=T^*BT+\eta_{\rm mem}^*\widetilde B\eta_{\rm mem},
\qquad
\mathcal Q_B(\psi)=
\mathcal Q_B(T\psi)+\mathcal Q_{\widetilde B}(\eta_{\rm mem}\psi).
\label{x_reciprocidad:nuc07:eq:17}
\end{equation}
\end{proposition}

\begin{proof}
Expand \(UJ\psi=JT\psi+\eta_{\rm mem}\psi\) in \(J^*U^*\widetilde BUJ=J^*\widetilde BJ\). Commutation with \(\Pi\) cancels the cross terms. Evaluating the operator identity and multiplying by \(\hbar\) gives the action balance.
\end{proof}

By iteration,

\begin{equation}
\mathcal Q_B(\psi)=\mathcal Q_B(T^n\psi)
+\sum_{j=0}^{n-1}\mathcal Q_{\widetilde B}
(\eta_{\rm mem}T^j\psi).
\label{x_reciprocidad:nuc07:eq:18}
\end{equation}

This sum records successive complements of the compressed dynamics. The complete dynamics has another exact formula, with \(z_n=U^nJ\psi\):

\begin{equation}
\mathcal Q_B(\psi)=
\mathcal Q_B(J^*z_n)+
\mathcal Q_{\widetilde B}((I-\Pi)z_n).
\label{x_reciprocidad:nuc07:eq:19}
\end{equation}

Without an additional property, \(T^n\) is not identified with \(J^*U^nJ\): the complete state can reabsorb memory. If \(B,\widetilde B\ge0\), all summands in \eqref{x_reciprocidad:nuc07:eq:18} are nonnegative; for signed charges, the equality remains, without this positivity interpretation.

If \(U\) is the propagator of an action such as \eqref{x_reciprocidad:nuc07:eq:13}, commutation of \(\widetilde B\) with its Hamiltonian supplies Noether conservation. The reduction condition involving \(\Pi\) specifies the readings that separate charge without interference terms.

Without imposing \eqref{x_reciprocidad:nuc07:eq:16}, the general identity preserves the additional terms:

\begin{equation}
\begin{aligned}
&\mathcal Q_B(T\psi)+\mathcal Q_{\widetilde B}(\eta_{\rm mem}\psi)
+2\hbar\operatorname{Re}
\langle JT\psi,\widetilde B\eta_{\rm mem}\psi\rangle\\
&=\mathcal Q_B(\psi)+
\hbar\langle\psi,J^*(U^*\widetilde BU-\widetilde B)J\psi\rangle.
\end{aligned}
\label{x_reciprocidad:nuc07:eq:20}
\end{equation}

The physical charge uses its reader, for example \(N_{108}\); a norm identity corresponds to \(B=I\). The sections \(\hbar_{\rm pre}\) and \(\hbar_{\rm ret}\) preserve the return reader of \cref{x_reciprocidad:iv:action:secciones,x_reciprocidad:iv:action:razon}; their difference and the memory complement have their respective domains of definition.

\section{Covariant balance of the total tensor}

The variational elimination of contorsion and the Noether identity of the covariant action, developed in Proposition~\ref{x_reciprocidad:vii:prop:fuente-metrica-balance}, provide

\begin{equation}
\mathcal T_{\mu\nu}
=T^{\rm phys,LC}_{\mu\nu}+U^{\rm phys,spin}_{\mu\nu},
\qquad
\mathring\nabla^\mu\mathcal T_{\mu\nu}=0.
\label{x_reciprocidad:nuc07:eq:21}
\end{equation}

Its domain is the realized HMT–MD chart, with the couplings held constant during the variation. For a Killing field \(\xi\),

\begin{equation}
j^\mu_\xi=\mathcal T^{\mu\nu}\xi_\nu,\qquad
\mathring\nabla_\mu j^\mu_\xi=0.
\label{x_reciprocidad:nuc07:eq:22}
\end{equation}

The proof uses \eqref{x_reciprocidad:nuc07:eq:21}, the symmetry of \(\mathcal T\), and the antisymmetry of \(\mathring\nabla_\mu\xi_\nu\). Upon integration, the difference between the charges of the hypersurfaces is the flux through the lateral boundary.

Conservation belongs to the total tensor; each component may exchange charge. Theorem~\ref{x_reciprocidad:vii:thm:intercambio-traza-residual} explicitly determines \(\mathcal J^{\rm tr}_\nu=\partial_\nu\operatorname{tr}_g(T^{\rm res})/4\) and the opposite balances of the trace and traceless part. This is recovered as a covariant instance of the joint-balance criterion, not as a new gravitational equation.
